% Standalone Fig. 2: per-value encoding lanes feeding a record memory map.
% The manuscript includes images/figure2-record-map.pdf exported from this source.
% Facts follow evidence/record-example.json: kernel-order scalars 0-2 are
% 3C12/3956/4178, with codes 000/111/111, escape selectors 0001/1111 (packed
% into byte F1), and raw patch (escape index 1, high byte 41). Low 16 KiB,
% codes 6 KiB, variable region 64 B = 48 + 1 + 3 + 4 + 8. Heights are not to
% scale. Coordinates are final-size points (252 pt = \columnwidth); y grows
% downward. Each lane's centre equals the centre of its record row.
% The manuscript's current Fig. 2 caption applies unchanged; at this height,
% replacing figure2-record-split.pdf keeps the four-page layout.
% Export: python3 paper/render_figures.py --figures 2 \
%   --source figures/record_layout_map.tex --output-stem figure2-record-map
\definecolor{rmInk}{HTML}{202C3A}%
\definecolor{rmMuted}{HTML}{6A788A}%
\definecolor{rmFrame}{HTML}{AAB8C7}%
\definecolor{rmPad}{HTML}{E9EDF2}%
\definecolor{rmBlue}{HTML}{246AC0}%
\definecolor{rmBlueTint}{HTML}{EDF4FC}%
\definecolor{rmTeal}{HTML}{00798C}%
\definecolor{rmTealTint}{HTML}{E6F5F7}%
\definecolor{rmOrange}{HTML}{C76821}%
\definecolor{rmOrangeTint}{HTML}{FFF3E6}%
\begin{tikzpicture}[x=1pt,y=-1pt,line cap=round,line join=round,
    every node/.style={inner sep=0pt,text=rmInk,
      font=\fontfamily{phv}\fontsize{6.8}{8}\selectfont},
    title/.style={anchor=base west,
      font=\fontfamily{phv}\fontseries{b}\fontsize{7.4}{8.8}\selectfont},
    small/.style={font=\fontfamily{phv}\fontsize{6.2}{7.2}\selectfont},
    lane/.style={anchor=base west,font=\fontfamily{phv}\fontsize{6.6}{7.8}\selectfont},
    mono/.style={font=\fontfamily{pcr}\fontseries{b}\fontsize{7}{8}\selectfont},
    chip/.style={mono,draw,fill=white,rounded corners=1.1pt,line width=.5pt,
      minimum height=8.6pt,inner xsep=2.3pt,text height=4.2pt,text depth=0pt},
    chain/.style={line width=.55pt,draw=rmOrange,-{Stealth[length=2.5pt,width=2.3pt]}}]
  \path[use as bounding box] (0,0) rectangle (252,115.5);

  \def\cA{60}\def\cB{97}\def\cC{134}% example columns
  \def\rIn{33.5}\def\rLow{46}\def\rCode{57}\def\rEsc{84.5}\def\rRaw{101}% lanes
  \def\recL{166}\def\recR{251.75}

  % #1 centre, #2 tint, #3 colour, #4 label: one lane that runs into the record.
  \def\rmLane#1#2#3#4{%
    \fill[#2,rounded corners=5pt] (.25,#1-5) rectangle (\recL+4,#1+5);
    \node[lane,text=#3] at (4,#1+2.3) {#4};
    \fill[#3] (\recL-6,#1-2.5) -- (\recL-1.2,#1) -- (\recL-6,#1+2.5) -- cycle;}
  % #1 centre, #2 high byte, #3 low byte: a solid two-byte BF16 value.
  \def\rmValue#1#2#3{%
    \begin{scope}
      \clip[rounded corners=1.1pt] (#1-15,\rIn-4.7) rectangle (#1+15,\rIn+4.7);
      \fill[rmTeal] (#1-15,\rIn-4.7) rectangle (#1,\rIn+4.7);
      \fill[rmBlue] (#1,\rIn-4.7) rectangle (#1+15,\rIn+4.7);
    \end{scope}
    \node[mono,text=white,text height=4.2pt,text depth=0pt] at (#1-7.5,\rIn) {#2};
    \node[mono,text=white,text height=4.2pt,text depth=0pt] at (#1+7.5,\rIn) {#3};}
  % #1 centre, #2 path, #3 table entry or patch, #4 recovered high byte.
  \def\rmPath#1#2#3#4{%
    \node[small,anchor=base,font=\fontfamily{phv}\fontseries{b}\fontsize{6.4}{7.4}\selectfont]
      at (#1,16.6) {#2};
    \node[small,anchor=base,text=rmMuted] at (#1,23.2)
      {#3\,$\rightarrow$\,{\fontfamily{pcr}\fontseries{b}\selectfont\color{rmTeal}#4}};}

  % --- Headers.
  \node[title] at (0,6.6) {Encoding individual BF16 values};
  \node[title] at (\recL,6.6) {KVFold record layout};
  \node[small,anchor=base west,text=rmMuted] at (\recL,16.6)
    {K record, 128\,\texttimes\,128 BF16};

  % --- Left: three example values, one per column, each with its decode path.
  \rmPath{\cA}{Direct}{P[0]}{3C}
  \rmPath{\cB}{Extended}{E[1]}{39}
  \rmPath{\cC}{Raw}{patch}{41}

  \rmLane{\rLow}{rmBlueTint}{rmBlue}{Low byte}
  \rmLane{\rCode}{rmTealTint}{rmTeal}{3-bit code}
  \rmLane{\rEsc}{rmOrangeTint}{rmOrange}{4-bit selector}
  \rmLane{\rRaw}{rmOrangeTint}{rmOrange}{Raw patch}

  \node[lane,text=rmMuted] at (4,\rIn+2.3) {BF16};
  \rmValue{\cA}{3C}{12}
  \rmValue{\cB}{39}{56}
  \rmValue{\cC}{41}{78}

  \foreach \c/\v in {\cA/12,\cB/56,\cC/78}
    \node[chip,draw=rmBlue,text=rmBlue] at (\c,\rLow) {\v};
  \node[chip,draw=rmTeal,text=rmTeal] at (\cA,\rCode) {000};
  \node[chip,draw=rmTeal,text=rmTeal] (b2) at (\cB,\rCode) {111};
  \node[chip,draw=rmTeal,text=rmTeal] (c2) at (\cC,\rCode) {111};
  \node[chip,draw=rmOrange,text=rmOrange] (b3) at (\cB,\rEsc) {0001};
  \node[chip,draw=rmOrange,text=rmOrange] (c3) at (\cC,\rEsc) {1111};
  \node[chip,draw=rmOrange,text=rmOrange] (c4) at (\cC,\rRaw) {(1,\,41)};
  \draw[chain] (b2.south) -- (b3.north);
  \draw[chain] (c2.south) -- (c3.north);
  \draw[chain] (c3.south) -- (c4.north);

  % --- Right: the record as a memory map. Spine colours match the lanes.
  \fill[rmBlueTint] (\recL,40.5) rectangle (\recR,51.5);
  \fill[rmTealTint] (\recL,51.5) rectangle (\recR,62.5);
  \fill[rmOrangeTint] (\recL,62.5) rectangle (\recR,115);
  \draw[rmFrame,line width=.5pt] (\recL,40.5) rectangle (\recR,62.5);
  \draw[rmFrame,line width=.5pt] (\recL+3,51.5) -- (\recR,51.5);
  \node[anchor=base west] at (\recL+7,\rLow+2.4) {Low bytes};
  \node[anchor=base east,text=rmMuted] at (\recR-4,\rLow+2.4) {16 KiB};
  \node[anchor=base west] at (\recL+7,\rCode+2.4) {3-bit codes};
  \node[anchor=base east,text=rmMuted] at (\recR-4,\rCode+2.4) {6 KiB};

  \node[anchor=base west,font=\fontfamily{phv}\fontseries{b}\fontsize{6.8}{8}\selectfont,
    text=rmOrange] at (\recL+7,68.9) {Variable region};
  \node[anchor=base east,text=rmMuted] at (\recR-4,68.9) {64 B};
  % #1 top, #2 bottom, #3 fill, #4 name, #5 size: one field of the variable region.
  \def\rmField#1#2#3#4#5{%
    \fill[#3] (\recL+6,#1) rectangle (\recR-3,#2);
    \node[lane] at (\recL+9,{(#1+#2)/2+2.3}) {#4};
    \node[lane,anchor=base east,text=rmMuted] at (\recR-6,{(#1+#2)/2+2.3}) {#5};}
  \rmField{72}{80}{white}{Header}{48 B}
  \rmField{80}{89}{white}{Selectors}{1 B}
  \rmField{89}{96.5}{rmPad}{Pad}{3 B}
  \rmField{96.5}{105.5}{white}{Raw}{4 B}
  \rmField{105.5}{112.5}{rmPad}{Pad}{8 B}
  \draw[rmOrange!45,line width=.4pt] (\recL+6,72) rectangle (\recR-3,112.5);
  \foreach \y in {80,89,96.5,105.5}
    \draw[rmOrange!45,line width=.4pt] (\recL+6,\y) -- (\recR-3,\y);
  \draw[rmOrange,line width=.55pt] (\recL,62.5) rectangle (\recR,115);
  % Spines are drawn last and extend half a stroke beyond the frames so that
  % no gray or orange outline shows along their outer edges.
  \fill[rmBlue] (\recL-.3,40.25) rectangle (\recL+3,51.5);
  \fill[rmTeal] (\recL-.3,51.5) rectangle (\recL+3,62.5);
  \fill[rmOrange] (\recL-.3,62.5) rectangle (\recL+3,115.3);
\end{tikzpicture}%
